\documentclass[11pt,a4paper]{article}
\usepackage[margin=1in]{geometry}
\usepackage{booktabs}
\usepackage{longtable}
\usepackage{graphicx}
\usepackage{hyperref}
\usepackage{listings}
\usepackage{xcolor}
\usepackage{enumitem}
\usepackage{caption}
\usepackage{amsmath}
\usepackage{amssymb}
\usepackage{tcolorbox}
\tcbuselibrary{breakable}

\hypersetup{
    colorlinks=true,
    linkcolor=blue,
    filecolor=magenta,
    urlcolor=cyan,
}

\tcbset{
    promptbox/.style={
        colback=blue!5!white,
        colframe=blue!75!black,
        title={#1},
        fonttitle=\bfseries,
        before skip=6pt,
        after skip=6pt,
        breakable,
    }
}

\title{Performance Report 20: Vulkan Engine \\
\large Water Rendering After Report 19 --- Measuring and Attacking the Intrinsic Cost}
\author{}
\date{\today}

\begin{document}

\maketitle

\begin{abstract}
Reports~17--19 audited the water pipeline from the outside in: the hybrid ray-traced paths (17, 18) and then the
frame-level machinery that survived the \textbf{Minimal} preset (19). Every one of report~19's twelve issues has now
landed (\texttt{3771c50} $\rightarrow$ \texttt{fa7ecdb}), so the redundant work is gone: one cull instead of two, no
tessellation stages when tessellation is off, a demand-gated back-face pass, single-attachment water passes when blur is
off, one UBO write, cached descriptors, cleared fast paths. Water is still slow. This report therefore turns inward and
\emph{measures} what is left, by compiling the shipped water modules and ablating them subsystem by subsystem.

The result is that the remaining cost is not redundancy --- it is the constant factor of the procedural surface itself.
The shipped non-RT water fragment module contains \textbf{46,320 SPIR-V instructions} and \textbf{3,073 inlined PCG hash
sites}, i.e. \textbf{24 four-dimensional Perlin FBM call sites}, of which \textbf{33,432 instructions are integer hash
ALU}. A wavy frame executes on the order of \textbf{75 four-dimensional Perlin evaluations per water pixel}
($\sim$9,600 PCG hashes, $\sim$8.6$\times10^4$ integer ALU ops); the \emph{default} configuration, which has waves off,
still executes 13 of them purely to wobble a screen-space refraction lookup. The non-tessellated water vertex shader is
\textbf{8,464 instructions} against \textbf{137} for the solid vertex shader, and it runs twice per frame (geometry $+$
back-face).

The report finds \textbf{14 issues} and pairs each with a copy-pasteable implementation prompt. The headline is a single
cross-cutting change: \emph{every} water noise evaluation is 4D Perlin (16 corners $\times$ 8 hashes) even though the
animation is already carried by domain translation; ablating the fourth axis alone cuts the water fragment module from
46,320 to \textbf{20,508} instructions ($-56\%$) and the hash ALU by $62.5\%$. Everything else --- cheap refraction
noise, a band-limited caustic curvature, evaluation-time distance LOD, occlusion rejection against the solid depth, a
smaller water target --- composes on top of it.
\end{abstract}

\tableofcontents
\newpage

% =====================================================================
\section{Executive Summary}
% =====================================================================

\subsection{What Report 19 Fixed, and Why That Was Not Enough}

Report~19's twelve items are all in the tree: C1 (non-tessellation pipeline family), C2 (caustics/specular/glitter
gates), C3 (demand-gated back-face pass, single cull), H4 (single-attachment variant), H5 (RT-off suppression), H6
(TES/VS gates and varying pruning), M7 (single UBO write $+$ cached descriptors), M8 (empty-water fast path), M9
(auxiliary sync cleanup), M10 (\texttt{glslc -O} and dead chains), L11 (dead targets), L12 (water quality tier). They
removed \emph{multiplicity} and \emph{dead work}: duplicated evaluations, unconditional passes, uncached CPU work.

What they did not touch is the \emph{constant factor}. Every surviving evaluation is a four-dimensional Perlin noise
built from a 16-corner hypercube whose every corner is an 8-hash PCG gradient. The water fragment stage performs 24 of
these FBM chains; the water vertex stage performs 4 more per vertex. Report~19's gates decide \emph{whether} those chains
run; they do not make them cheaper. With the default layer (\texttt{enableWaves = false}) the gated work is off, and the
entire remaining procedural cost is one function --- \texttt{waterRefractionNoise} --- which costs 13 four-dimensional
Perlin evaluations per pixel to offset a screen-space UV.

\subsection{Measured Module Budget}

All numbers below come from \texttt{spirv-dis} on the shipped modules (\texttt{bin/shaders/main\_water.frag.spv},
\texttt{main\_water.tesc.spv}, \texttt{main\_water.tese.spv}, \texttt{main\_water\_no\_tess.vert.spv}) and from
controlled ablations of the same sources compiled with the Makefile's exact flags
(\texttt{glslc --target-env=vulkan1.3 -O -DWATER\_MODE=1}). The ablation protocol is in \S\ref{sec:method}.

\begin{table}[h]
\centering
\small
\begin{tabular}{@{}lrr@{}}
\toprule
\textbf{Shipped water module} & \textbf{SPIR-V instr.} & \textbf{Inlined PCG hash sites} \\
\midrule
\texttt{main\_water.frag} (non-RT water fragment) & 46,320 & 3,073 \\
\texttt{main\_water.tese} (tessellated water vertex) & 8,604 & 513 \\
\texttt{main\_water\_no\_tess.vert} (water vertex, tess off) & 8,464 & 513 \\
\texttt{main\_water.tesc} (water tessellation control) & 7,682 & 513 \\
\texttt{main\_water.vert} (plain water VS, tess path) & 137 & 0 \\
\texttt{main.frag} (solid fragment, for scale) & --- & 0 \\
\bottomrule
\end{tabular}
\caption{Static module sizes. A ``PCG hash site'' is one inlined \texttt{pcgHash} (\texttt{perlin.glsl:9}); 128 of them
make one four-dimensional Perlin evaluation (16 corners $\times$ \texttt{hash44} = 8 PCG each).}
\end{table}

\begin{table}[h]
\centering
\small
\begin{tabular}{@{}lrrl@{}}
\toprule
\textbf{Ablation of \texttt{main\_water.frag}} & \textbf{PCG sites} & $\Delta$ & \textbf{FBM sites removed} \\
\midrule
Baseline & 3,073 & --- & --- \\
Remove \texttt{waterRefractionNoise} & 2,561 & $-512$ & 4 \\
Remove specular $+$ glitter noise & 2,689 & $-384$ & 3 \\
Remove the caustic field pair & 2,049 & $-1{,}024$ & 8 \\
Wave field reduced to one ridged train & 1,409 & $-1{,}664$ & 13 \\
Perlin 4D $\rightarrow$ 3D (drop the time axis) & 1,153 & $-1{,}920$ & $-62.5\%$ hash ALU, 20,508 instr. \\
\bottomrule
\end{tabular}
\caption{Attribution. The 24 FBM call sites split exactly: 17 wave-field (5 shading $+$ 8 caustics $+$ 4 debug), 4
refraction noise, 3 specular/glitter.}
\end{table}

\subsection{Per-Pixel Dynamic Cost}

With the default spectrum (\texttt{noiseOctaves = 4}, \texttt{noiseLacunarity = 4.0}) each FBM site executes
\texttt{octaves} four-dimensional Perlin evaluations:

\begin{table}[h]
\centering
\small
\begin{tabular}{@{}lrcr@{}}
\toprule
\textbf{Consumer} & \textbf{FBM sites} & \textbf{Perlin4D / pixel} & \textbf{Gate} \\
\midrule
Shading wave field (\texttt{withFoam=true}) & 5 & 20 & \texttt{waveToggles.x} \\
Caustic curvature (2$\times$ \texttt{waterWaveSample}) & 8 & 32 & thickness $>0.05$, waves, intensity \\
Refraction screen offset & 4 & 13 & \texttt{enableRefraction} \\
Specular $+$ glitter & 3 & 10 & waves $+$ intensities \\
Debug displacement view & 4 & 16 & debug mode only \\
\midrule
\textbf{Default preset} (waves off, refraction on) & --- & \textbf{13} & --- \\
\textbf{Waves on, no caustics} & --- & \textbf{43} & --- \\
\textbf{Waves on, thick water (caustics)} & --- & \textbf{75} & --- \\
\bottomrule
\end{tabular}
\caption{Dynamic four-dimensional Perlin evaluations per shaded water pixel. One evaluation is $\sim$128 PCG hashes
$\approx$ 1.1--1.4$\times10^3$ ALU ops.}
\end{table}

\begin{itemize}
    \item \textbf{C1 --- The noise is 4D where it only needs to be 3D.} \texttt{hash44} (\texttt{perlin.glsl:34}) costs 8
    PCG calls per corner and \texttt{perlinNoise4D} (\texttt{:70}) samples 16 corners, so every Perlin evaluation costs
    128 PCG hashes. But the water field already animates by \emph{domain translation}: the chop, mask and foam FBMs
    subtract \texttt{shoreDrift}/\texttt{foamDrift} (\texttt{water\_noise.glsl:225,239,361}) and both ridged trains fold
    $-c\cdot t$ into their along-axis coordinate (\texttt{:76--93}). The fourth axis is a second, redundant animation
    mechanism. Ablation: 3,073 $\rightarrow$ 1,153 PCG sites, 46,320 $\rightarrow$ 20,508 instructions.
    \item \textbf{C2 --- \texttt{waterRefractionNoise} is 13 four-dimensional Perlin evaluations for a 2-channel
    screen-space offset.} It is the \emph{entire} procedural cost of the default configuration, it is not distance-faded,
    and with \texttt{noiseLacunarity = 4.0} its finest octave has a 6\,cm wavelength --- sub-pixel aliasing paid at full
    ALU cost (\texttt{water\_noise.glsl:398--410}, call site \texttt{water\_surface.glsl:714--723}).
    \item \textbf{C3 --- Caustics cost more than the surface they decorate.} The curvature term is a central difference
    of the \emph{full} field (32 of the 75 evaluations), i.e. two complete wave fields per pixel, to obtain one scalar
    $d^2h/du^2$ (\texttt{water\_surface.glsl:1458--1462}).
    \item \textbf{C4 --- The distance LOD is applied to the result, not to the evaluation.} \texttt{water\_surface.glsl:688--689}
    fades the wave normal to flat over 200--1000\,m, but the field is evaluated \emph{before} that at full octave count
    for every water pixel --- including every pixel whose result is then discarded. Refraction noise, specular and
    glitter have no distance gate at all.
    \item \textbf{C5 --- Terrain-occluded water is shaded in full and thrown away at composite time.} The water pass has
    no solid-depth rejection; the occlusion test lives in \texttt{postprocess.frag:147--150}. The water task already waits
    \texttt{tlSolid} (\texttt{MyApp.cpp:2242}) and \texttt{solidSceneDepthTex} is already bound in set~2, so one fetch and
    one \texttt{discard} at the top of the water fragment stage removes the whole shader for those pixels. (The stale
    comment at \texttt{WaterRenderer.cpp:1250} still claims ``forward-pass depth-test handles solid occlusion''.)
\end{itemize}

\textbf{Recommended order:} C1 (drop the fourth axis) $\rightarrow$ C2 (cheap refraction noise) $\rightarrow$ C4
(evaluation-time LOD) $\rightarrow$ C5 (occlusion rejection) $\rightarrow$ C3 (band-limited caustics) $\rightarrow$ H6
(band-limit the vertex stage, undisplaced back faces) $\rightarrow$ H7 (cheap TCS modulation) $\rightarrow$ H8/H9
(target format and resolution) $\rightarrow$ M10--M13 $\rightarrow$ L14.

\begin{table}[h]
\centering
\small
\begin{tabular}{@{}ll@{}}
\toprule
\textbf{Remaining water cost} & \textbf{Covered by} \\
\midrule
24$\times$ 4D-Perlin FBM chains, 33k integer ALU instructions & C1 \\
13 Perlin4D for the refraction screen offset (default preset) & C2, L14 \\
32 Perlin4D for the caustic curvature & C3 \\
Full-cost evaluation at distances where the result is discarded & C4 \\
Full shading of terrain-occluded water pixels & C5, H9 \\
16 Perlin4D per water vertex, twice per frame & H6 \\
3 full field evaluations per patch for tessellation modulation & H7 \\
RGBA32F full-res water target (16\,B/px write $+$ read $+$ clear) & H8 \\
Separate full-res water pass instead of water-in-main & H9 \\
10 Perlin4D for specular/glitter with no lobe early-out & M10 \\
Vertex and fragment evaluation of the same field & M11 \\
Full-res \texttt{waterGeomDepth} fetch for every screen pixel & M12 \\
Full pass for water that is resident but entirely frustum-culled & M13 \\
One shared noise spectrum for five consumers & L14 \\
\bottomrule
\end{tabular}
\end{table}

% =====================================================================
\section{Measurement Method}
\label{sec:method}
% =====================================================================

Static counts are taken from the shipped SPIR-V:

{\small
\begin{verbatim}
spirv-dis bin/shaders/main_water.frag.spv \
  | grep -oE "uint_747796405" | wc -l      # PCG hash sites
spirv-dis bin/shaders/main_water.frag.spv \
  | grep -oE "Op[A-Za-z]+" | sort | uniq -c
\end{verbatim}
}

\texttt{pcgHash} (\texttt{perlin.glsl:9--14}) multiplies by the literal \texttt{747796405u} exactly once, so counting that
literal counts inlined hash evaluations. Each \texttt{hash44} performs 8 of them (4 chained over the lattice coordinates
plus 4 for the output components); each four-dimensional Perlin evaluation performs 16 \texttt{hash44}. Hence
$3{,}073 / 128 \approx 24$ FBM call sites in the fragment module.

Ablations recompile \texttt{shaders/main.frag} from a scratch copy of \texttt{shaders/} with the Makefile's water flags
(\texttt{glslc --target-env=vulkan1.3 -Iincludes -O -DWATER\_MODE=1}) and one subsystem stubbed at a time. Because the
octave loops are dynamic (\texttt{octaves} is a UBO field), the compiler does \emph{not} unroll them: one FBM call site
in the source is one Perlin body in the module, which is why static site counts map cleanly onto dynamic cost.

Instruction mix of the shipped fragment module (dominant opcodes):

\begin{center}
\small
\begin{tabular}{@{}lr@{}}
\toprule
Opcode & count \\
\midrule
\texttt{OpShiftRightLogical} & 10,752 \\
\texttt{OpIAdd} & 7,704 \\
\texttt{OpBitwiseXor} & 7,296 \\
\texttt{OpIMul} & 6,144 \\
\texttt{OpFSub} & 2,574 \\
\texttt{OpBitcast} & 1,921 \\
\texttt{OpBitwiseOr} & 1,536 \\
\texttt{OpExtInst} (of which \texttt{Pow} $\times$28) & 964 \\
\texttt{OpFMix} & 531 \\
\bottomrule
\end{tabular}
\end{center}

Integer hash ALU (shift $+$ iadd $+$ xor $+$ imul $+$ or) is \textbf{33,432} of the 46,320 instructions: \textbf{72\% of
the water fragment module is spent computing random numbers}. That single fact drives the priority order below.

% =====================================================================
\section{Cost Model}
\label{sec:cost}
% =====================================================================

Let $P$ be the four-dimensional Perlin evaluation ($\approx 128$ PCG hashes $\approx 1.3\times10^3$ ALU ops). With
$O = \texttt{noiseOctaves}$:

\begin{equation}
C_{\mathrm{water}} = \underbrace{N_{\mathrm{px}} \cdot O \cdot (n_{\mathrm{field}} + n_{\mathrm{caustic}} + n_{\mathrm{refr}} + n_{\mathrm{spec}})\,P}_{\text{per-pixel procedural}}
\;+\; \underbrace{N_{\mathrm{vtx}} \cdot O \cdot n_{\mathrm{vtx}}\,P}_{\text{per-vertex, $\times 2$ passes}}
\;+\; \underbrace{N_{\mathrm{px}} \cdot C_{\mathrm{bw}}}_{\text{target bandwidth}}
\;+\; C_{\mathrm{fixed}}
\end{equation}

with the default $O = 4$ and

\[
n_{\mathrm{field}} = 5,\quad n_{\mathrm{caustic}} = 8,\quad n_{\mathrm{refr}} = 4,\quad n_{\mathrm{spec}} = 3,\quad
n_{\mathrm{vtx}} = 4 \ \text{(plus 3 more per patch in the TCS)}.
\]

At 1080p with 40\% water coverage and 60\,Hz that is $\sim 5\times10^7$ water-pixel invocations per second; at
$13P \approx 1.7\times10^4$ dynamically executed instructions each (default preset) the water pixel shader alone issues
$\sim 8\times10^{11}$ instructions/second, and at $75P$ (waves $+$ caustics) roughly $5\times10^{12}$ --- i.e. the water
pass by itself can exceed the whole-device issue budget of a mid-range GPU, before any rasterisation, texturing, or
bandwidth. This is consistent with the observed behaviour and explains why removing passes (report~19) improved things
without fixing them.

The three levers, in descending leverage:

\begin{enumerate}
    \item \textbf{Reduce $P$} --- the per-evaluation constant factor (C1: 4D $\rightarrow$ 3D, $-62.5\%$ hash ALU).
    \item \textbf{Reduce the $n$'s} --- do not evaluate what is not used (C2, C3, C4, M10, H6, H7).
    \item \textbf{Reduce $N$} --- shade fewer pixels (C5, H9) and move fewer bytes (H8).
\end{enumerate}

% =====================================================================
\section{Critical Issues}
% =====================================================================

\subsection{C1: Every Water Noise Evaluation Is Four-Dimensional, But the Animation Is Already Domain Translation}

\textbf{Location:} \texttt{shaders/includes/perlin.glsl:9--14} (\texttt{pcgHash}), \texttt{:34--44} (\texttt{hash44}, 8 PCG
per corner), \texttt{:70--113} (\texttt{perlinNoise4D}, 16 corners), \texttt{:156--236} (\texttt{perlinNoise4DGrad}),
\texttt{:133--147}/\texttt{:244--259} (\texttt{fbm}/\texttt{fbmGrad4D}); consumers
\texttt{shaders/includes/water\_noise.glsl:11--16} (\texttt{waterFbmNoise}), \texttt{:19--26}
(\texttt{waterFbmNoiseGrad}), \texttt{:35--67} (\texttt{waterRidgedFbmGrad}).

\begin{itemize}
    \item A four-dimensional Perlin evaluation costs \textbf{16 corner hashes}, and \texttt{hash44} costs \textbf{8 PCG
    calls} (four chained over the integer lattice coordinates, four for the output components). One evaluation is
    therefore \textbf{128 PCG hashes} $\approx$ 700 integer ALU ops before the interpolation tree. The shipped water
    fragment module contains 3,073 of them.
    \item The fourth axis is the \emph{time} axis, and it is redundant: every consumer already animates by translating
    its sample point. The chop, the calm mask and the foam texture subtract a drift term
    (\texttt{water\_noise.glsl:225, 239, 361} --- \texttt{xyz - shoreDrift} / \texttt{xyz - foamDrift}); both ridged
    trains fold the phase directly into their along-axis coordinate,
    \texttt{along = dot(pos.xz, moveDir) + phaseOffset - speed * speedFactor * time} (\texttt{:76--93}). Passing
    \texttt{time} as the fourth Perlin axis is a \emph{second}, independent animation mechanism on top of that.
    \item Measured effect of removing the fourth axis (8 corners $\times$ \texttt{hash33} = 6 PCG $=48$ PCG per
    evaluation): \textbf{3,073 $\rightarrow$ 1,153 PCG sites ($-62.5\%$)} and \textbf{46,320 $\rightarrow$ 20,508 SPIR-V
    instructions ($-55.7\%$)}. This is a single, cross-cutting change that makes every other item cheaper too.
    \item The only consumer with no drift of its own is \texttt{waterRefractionNoise}
    (\texttt{water\_noise.glsl:398--410}), which passes a static \texttt{xyz}; it needs an explicit scroll term
    (see C2).
\end{itemize}

\textbf{Technique:} add 3D gradient/noise entry points next to the 4D ones and switch the water FBM helpers to them,
compensating the loss of the time axis with the drift each consumer already applies (and with an explicit scroll where
there is none). Keep the 4D paths for non-water users so nothing else regresses.

\begin{tcolorbox}[promptbox={Prompt: Drop the Fourth Noise Axis from the Water Field}]
\textbf{Goal:} One water noise evaluation costs 48 PCG hashes instead of 128; the water fragment module drops below
21{,}000 instructions; the animated look is unchanged in character.

\textbf{Steps:}
\begin{enumerate}
    \item Add \texttt{vec4 perlinNoise3DGrad(vec3 p)} to \texttt{perlin.glsl} beside \texttt{perlinNoise4DGrad}
    (\texttt{:156}): 8 corners of \texttt{hash33}, same Hermite/$30f^2(1-f)^2$ derivative bookkeeping, returning
    \texttt{vec4(value, d/dx, d/dy, d/dz)}. Mirror it with \texttt{vec4 fbmGrad3D(vec3 p, ...)} beside
    \texttt{fbmGrad4D} (\texttt{:244}). Keep the 4D versions for the solid/vegetation users.
    \item Switch \texttt{waterFbmNoise} / \texttt{waterFbmNoiseGrad} (\texttt{water\_noise.glsl:11--26}) to the 3D
    entry points, feeding the existing \texttt{time} argument into the sample position as an explicit drift
    (\texttt{xyz + driftDir * time * timeScale}) where the caller does not already subtract one --- the chop/mask/foam
    call sites already pass \texttt{xyz - shoreDrift}, so bump their drift magnitude to restore the previous apparent
    motion speed; \texttt{waterRefractionNoise} needs a new scroll term (C2).
    \item Switch \texttt{waterRidgedFbmGrad} (\texttt{:35--67}) to \texttt{perlinNoise3DGrad} on
    \texttt{vec3(q * freq, slice + phaseTime)}: the along-coordinate already carries $-c\,t$, and \texttt{slice}
    (0.0 / 17.0 for the two trains) becomes a constant offset on the third axis, preserving the decorrelation between
    trains. Verify the ridged crest geometry is unchanged in character (A/B screenshots at two times).
    \item Keep \texttt{hash44}/\texttt{perlinNoise4D} in the file: they are still used elsewhere and deleting them
    would move the solid path.
    \item Re-measure with the \S\ref{sec:method} protocol: expect 1,153 PCG sites and $\sim$20.5k instructions for
    \texttt{main\_water.frag}, and a proportional drop in \texttt{main\_water.tese} /
    \texttt{main\_water\_no\_tess.vert}.
\end{enumerate}

\textbf{Files to modify:}
\begin{itemize}
    \item \texttt{shaders/includes/perlin.glsl} (new 3D gradient entry points)
    \item \texttt{shaders/includes/water\_noise.glsl} (FBM helpers, ridged FBM, refraction-noise scroll)
    \item \texttt{Makefile} only if new variant rules are needed (none expected)
\end{itemize}

\textbf{Validation:} \texttt{spirv-dis} PCG-site count drops from 3,073 to $\sim$1,153 and instruction count to
$\sim$20.5k; screenshot A/B at $t$ and $t+\Delta$ shows the same surface character and the same apparent wave speed
(no frozen or boiling-in-place artefacts); the solid/vegetation passes are byte-identical in output; GPU timestamps
14/15 drop measurably on water-dominated views; validation clean, and GPU-assisted validation clean via
\texttt{VULKAN\_GPU\_ASSISTED=1 make run-debug}.
\end{tcolorbox}

\subsection{C2: \texttt{waterRefractionNoise} Spends 13 Four-Dimensional Perlin Evaluations on a Two-Channel Screen Offset}

\textbf{Location:} \texttt{shaders/includes/water\_noise.glsl:398--410}; call site
\texttt{shaders/includes/water\_surface.glsl:714--723}; consumers \texttt{:728--732} (screen-space offset) and
\texttt{:846--850} (Snell-ray perturbation); defaults \texttt{utils/WaterParams.hpp:37,39}
(\texttt{noiseOctaves = 4}, \texttt{noiseLacunarity = 4.0}).

\begin{itemize}
    \item The helper evaluates four separate FBM chains --- \texttt{noise1} (4 octaves), \texttt{noise2} (3),
    \texttt{noise3} (2), and a second 4-octave chain for \texttt{nY} --- i.e. \textbf{13 four-dimensional Perlin
    evaluations} to return a \texttt{vec2}. Measured ablation: 512 of the 3,073 PCG sites (17\% of the module's hash
    ALU) live in this one function.
    \item It is the \textbf{entire} procedural cost of the default configuration. Layer~0 ships with
    \texttt{enableWaves = false}, so the wave field early-outs, caustics are gated off, and specular/glitter take the
    flat branch (\texttt{water\_surface.glsl:1191--1194}). What is left of the water pixel shader is 13 Perlin4D
    evaluations plus a handful of texture fetches.
    \item Its product is a \emph{screen-space} quantity: \texttt{refractionOffset = refractionNoise *
    refractionStrength} either offsets \texttt{screenUV} for the raster bottom lookup (\texttt{:972--976}) or tilts the
    Snell direction (\texttt{:846--850}). It is not distance-faded, so water 3\,km away pays the same 13 evaluations for
    an offset whose screen footprint is a fraction of a pixel.
    \item With \texttt{noiseLacunarity = 4.0} and 4 octaves the finest octave is at $64\times$ the base frequency
    (base wavelength 4\,m $\Rightarrow$ 6.25\,cm features). Those octaves are invisible or aliased on screen at almost
    any distance, and they cost exactly as much as the octave that produces the visible motion.
\end{itemize}

\textbf{Technique:} replace the four chained FBMs with one cheap, band-limited, two-channel source: a single
2--3-octave 3D/2D noise with a scroll term, or a pre-baked RG noise texture sampled in world-XZ with a time scroll (the
engine already has a compute noise generator in \texttt{services/TextureMixer.cpp}). Gate it on
\texttt{refractionStrength > 0} and fade it out with camera distance and with the screen-space footprint of the offset.

\begin{tcolorbox}[promptbox={Prompt: Make the Refraction Distortion Cheap, Band-Limited and Distance-Aware}]
\textbf{Goal:} The refraction offset costs at most one low-octave noise evaluation (or one texture fetch) and vanishes
where it cannot be seen.

\textbf{Steps:}
\begin{enumerate}
    \item Early-out at the call site (\texttt{water\_surface.glsl:714}) when
    \texttt{refractionStrength <= 0.0} (and, for the raster path, when the layer has no bottom to sample): today the
    gate is only \texttt{enableRefraction || dbgMode}.
    \item Replace the body of \texttt{waterRefractionNoise} with a single 2-octave 3D FBM sampled at
    \texttt{xyz + drift}, returning two decorrelated channels from the same evaluation (e.g. two different
    \texttt{offset} seeds through the cheap 3D helper, or one \texttt{fbmGrad3D} plus a rotated re-sample). Keep the
    returned range and sign conventions so \texttt{refractionStrength} keeps its current meaning; document the look
    change (fewer, larger features).
    \item Alternative if a texture is preferred: add a 64$\times$64 RG16F (or RGBA8) tiling noise generated once at
    startup with the existing \texttt{perlin\_noise.comp} machinery, sample it at
    \texttt{fragPos.xz * scale + scroll * t}, and keep the two channels. Validate that tiling is invisible at the
    authored \texttt{noisePeriod}.
    \item Add a distance/footprint fade: compute the world-space size of a pixel at the fragment
    (\texttt{fwidth(fragPosWorld)} or the depth-derived scale) and scale the offset --- and skip the noise entirely ---
    once that footprint exceeds the finest retained octave's wavelength. This removes the cost on distant water, which
    is where most water pixels are.
    \item Revisit the defaults: \texttt{noiseLacunarity = 4.0} with 4 octaves is a $64\times$ frequency spread. Either
    lower the default to $\sim$2.0--2.5 or split the spectrum per consumer (L14) so the refraction offset never pays for
    octaves it cannot show.
\end{enumerate}

\textbf{Files to modify:}
\begin{itemize}
    \item \texttt{shaders/includes/water\_noise.glsl} (\texttt{waterRefractionNoise}, \texttt{:398--410})
    \item \texttt{shaders/includes/water\_surface.glsl} (call site \texttt{:714--723}, distance fade)
    \item \texttt{utils/WaterParams.hpp} / \texttt{vulkan/ubo/WaterParamsGPU.hpp} (defaults / new packed fields, if any)
\end{itemize}

\textbf{Validation:} PCG-site count drops by $\sim$512 (or by the full 4-site cost once C1 lands); a camera looking at
distant water shows no refraction noise cost in the shader counters while near water is unchanged; screenshot A/B of
still shallow water is within tolerance and the animated distortion still reads as water; validation clean, GPU-assisted
validation clean.
\end{tcolorbox}

\subsection{C3: The Caustic Curvature Evaluates the Full Wave Field Twice --- More Than the Surface Itself}

\textbf{Location:} \texttt{shaders/includes/water\_surface.glsl:1416--1470}, in particular the two
\texttt{waterWaveSample} calls at \texttt{:1458--1462}; field implementation
\texttt{shaders/includes/water\_noise.glsl:118--384}.

\begin{itemize}
    \item The caustic term needs one scalar, $d^2h/du^2$ along the sun azimuth. It obtains it as a central difference of
    the analytic gradient, which means \textbf{two complete evaluations of the wave field} at
    \texttt{entry $\pm$ sunHat * ec}. Each field evaluation is 4 FBM chains (chop, mask, train 1, train 2), so the
    curvature costs \textbf{8 FBM sites = 32 four-dimensional Perlin evaluations per pixel} --- 43\% of the 75 evaluations
    in a wavy, thick-water frame, and more than the 20 the shading normal costs. Measured ablation: removing the block
    removes 1,024 of the 3,073 PCG sites.
    \item The curvature is a \emph{second-derivative} quantity. It does not need the amplitude envelope, the calm mask,
    the foam train, or the second ridged train to be resolved at full fidelity: the caustic pattern is dominated by the
    primary crest geometry.
    \item The report~19 gates (\texttt{causticIntensity > 0.001 \&\& waterThickness > 0.05 \&\& waveToggles.x > 0.5},
    \texttt{water\_surface.glsl:1428--1429}) correctly prevent the cost on flat/thin water, but wherever the effect is
    actually visible it is the single largest term in the shader.
\end{itemize}

\textbf{Technique:} compute the curvature from a reduced field and/or from a one-sided stencil. Two cheap options:
(a) reuse the gradient already returned by the shading field at the shaded point and take a single one-sided difference
(1 field evaluation instead of 2); (b) add a dedicated \texttt{waterWaveCurvature()} that evaluates \emph{only} the
primary ridged train with a three-point stencil, skipping the chop, the mask, the second train and foam.

\begin{tcolorbox}[promptbox={Prompt: Band-Limit the Caustic Curvature to One Reduced Field Evaluation}]
\textbf{Goal:} The caustic term costs at most 4--8 four-dimensional Perlin evaluations instead of 32, with the same
caustic pattern at the authored softness.

\textbf{Steps:}
\begin{enumerate}
    \item Add \texttt{float waterWaveCurvature(vec3 pos, float time, float depth, float amp, vec2 shoreDir, vec3 u,
    float h, WaterParamsGPU wp)} to \texttt{water\_noise.glsl}: evaluate \emph{only} the primary ridged train
    (\texttt{waterRidgedFbmGrad} on \texttt{coord1}) at \texttt{pos}, \texttt{pos + u*h}, \texttt{pos - u*h} and return
    the second difference. Reuse \texttt{waterShoreWaveCoord} so the movement is identical to the surface. Cost: 3
    single-train evaluations $\approx$ 12 Perlin4D (4 with C1's 3D noise).
    \item Alternatively (cheaper, slightly different): reuse the analytic gradient \texttt{waveField.grad} already
    computed for the shading normal at \texttt{fragBasePos} and take one one-sided difference against a single field
    evaluation at the entry point. Document which approximation is chosen and why.
    \item Replace the \texttt{waveP}/\texttt{waveM} pair at \texttt{water\_surface.glsl:1458--1462} with the chosen
    helper; keep \texttt{ec} (the quarter-wavelength stencil from the finest retained octave) exactly as it is so the
    caustic band-limiting is unchanged.
    \item Preserve the \texttt{causticDebugMode} escape hatch: the debug view must still be able to force the accurate
    path (or accept the reduced one --- document the choice).
    \item Re-measure: expect the caustics block to fall from 1,024 PCG sites to $\sim$256--384.
\end{enumerate}

\textbf{Files to modify:}
\begin{itemize}
    \item \texttt{shaders/includes/water\_noise.glsl} (new \texttt{waterWaveCurvature}, or a
    \texttt{waterWaveFieldReduced} variant)
    \item \texttt{shaders/includes/water\_surface.glsl} (caustic block \texttt{:1416--1470})
\end{itemize}

\textbf{Validation:} \texttt{DEBUG\_MODE\_CAUSTICS} screenshots are within tolerance of the baseline at the default
softness (compare the fold pattern's spatial frequency and contrast, not exact pixel values); the caustic animation
still rides the waves; PCG-site count drops by $\geq$640; validation clean, GPU-assisted validation clean.
\end{tcolorbox}

\subsection{C4: The Distance LOD Is Applied to the Wave Field's Result, Not to Its Evaluation}

\textbf{Location:} \texttt{shaders/includes/water\_surface.glsl:665--667} (field evaluation),
\texttt{:686--694} (distance fade), \texttt{:714--723} (refraction noise), \texttt{:1174--1190} (specular/glitter),
\texttt{:1416--1470} (caustics).

\begin{itemize}
    \item The water fragment stage evaluates the wave field first (\texttt{:665}) and only afterwards computes
    \texttt{detail = 1.0 - smoothstep(200.0, 1000.0, camDist)} and mixes the resulting normal toward the flat base
    normal (\texttt{:688--694}). Every water pixel beyond 1\,km therefore pays the \textbf{full 20 four-dimensional
    Perlin evaluations} for a normal that is then discarded entirely, and every pixel in the 200--1000\,m band pays full
    price for a partially used result.
    \item Refraction noise (13), specular and glitter (10) and caustics (32) have \emph{no} distance gate whatsoever, so
    distant water pays 75 evaluations per pixel while producing a nearly flat, sky-mirroring surface.
    \item Distant water is not a small number of pixels: in a terrain view the water body recedes to the horizon, so the
    far band routinely covers more screen area than the near band. This is the cheapest large win in the report after C1.
\end{itemize}

\textbf{Technique:} hoist the camera distance and the pixel footprint above the field evaluation and use them to (a) skip
the field entirely past the fade distance, (b) reduce the octave count inside the fade band, and (c) gate the refraction,
specular, glitter and caustic blocks on the same quantity.

\begin{tcolorbox}[promptbox={Prompt: Evaluate at the LOD the Pixel Can Actually Resolve}]
\textbf{Goal:} Water beyond the detail fade distance performs zero wave/specular/glitter/caustic noise work; water in the
fade band performs a proportionally reduced octave count.

\textbf{Steps:}
\begin{enumerate}
    \item Move the \texttt{camDist}/\texttt{detail} computation to before the field call at
    \texttt{water\_surface.glsl:665} and expose an integer \texttt{octBudget} (e.g. \texttt{4} near, \texttt{2} mid,
    \texttt{0} beyond the fade) plus a hard \texttt{skipField} boolean. Pass \texttt{octBudget} into
    \texttt{waterWaveField} (a new trailing parameter) and clamp \texttt{octavesN} with it.
    \item When \texttt{skipField} is true, skip the field call entirely and use \texttt{flatN} as the normal; also skip
    the refraction noise, the caustic block and the specular/glitter FBMs (keep the analytic highlight if it is still
    wanted at range --- it costs one \texttt{pow}).
    \item Prefer a \emph{footprint}-based criterion over raw distance where possible: the world-space extent of a pixel
    (\texttt{fwidth(fragPosWorld)} or the depth-derived scale) is what actually band-limits the surface, and it
    automatically handles both grazing far water and near-field zoom.
    \item Apply the same budget to the vertex stage (see H6): the TES/no-tess VS currently evaluates the field at full
    octave count for every vertex regardless of distance.
    \item Keep the fade band's endpoints (200/1000\,m) as-is in the first iteration so the look is unchanged; A/B before
    moving them.
\end{enumerate}

\textbf{Files to modify:}
\begin{itemize}
    \item \texttt{shaders/includes/water\_surface.glsl} (hoist LOD, gate the blocks)
    \item \texttt{shaders/includes/water\_noise.glsl} (octave-budget parameter on \texttt{waterWaveField})
    \item \texttt{shaders/main.vert} / \texttt{shaders/includes/water\_tese.glsl} (vertex-side budget, H6)
\end{itemize}

\textbf{Validation:} A view dominated by distant water shows the wave/refraction/specular blocks skipped in a RenderDoc
capture (use the existing debug views or a temporary counter) while near water is pixel-identical; no popping appears at
the 200\,m / 1000\,m boundaries (fade the octave count, never the amplitude); GPU time on horizon-heavy views drops
measurably; validation clean.
\end{tcolorbox}

\subsection{C5: Terrain-Occluded Water Is Shaded in Full and Discarded Only at Composite Time}

\textbf{Location:} \texttt{shaders/includes/water\_surface.glsl:606--610} (the note recording that the
solid-occlusion discard was removed), \texttt{shaders/postprocess.frag:146--151} (the composite-side occlusion test),
\texttt{vulkan/renderer/WaterRenderer.cpp:1246--1252} (depth attachment cleared to 1.0, with the stale comment
``Forward-pass depth-test handles solid occlusion''), \texttt{MyApp.cpp:2242} and \texttt{:2678} (the water task waits
\texttt{tlSolid}).

\begin{itemize}
    \item The water geometry pass has its own depth target, cleared to 1.0, and writes only water. Nothing in it
    consults the solid depth, so a water surface hidden behind a ridge, a bank or an overhang is rasterised and shaded
    with the \emph{full} water shader --- up to 75 four-dimensional Perlin evaluations per pixel --- and is then thrown
    away by \texttt{postprocess.frag:148} (\texttt{if (waterGeomDepth < 1.0 \&\& obstacleDepth < waterGeomDepth)
    waterAlpha = 0.0;}).
    \item This is not a correctness requirement: the water task already waits on \texttt{tlSolid}, and
    \texttt{solidSceneDepthTex} is already bound in the water set~2 and sampled elsewhere in the same shader
    (\texttt{water\_surface.glsl:975, 922, 1091}). One fetch at the top of the water fragment stage is enough.
    \item The pixel cost is ``free'' only in the sense that a discard still rasterises; what it saves is the entire
    procedural shading stack, which is 72\% of the module.
    \item The comment at \texttt{WaterRenderer.cpp:1250} is actively misleading and should be corrected in the same
    change.
\end{itemize}

\textbf{Technique:} restore an explicit early rejection against the solid depth inside the water fragment stage, placed
before the wave field; make it conditional on the solid depth actually being available (it is, whenever the water task
runs after \texttt{tlSolid}, which is the only ordering in use).

\begin{tcolorbox}[promptbox={Prompt: Reject Solid-Occluded Water Fragments Before Shading}]
\textbf{Goal:} Water hidden behind terrain costs one depth fetch and a discard, not the full water shader.

\textbf{Steps:}
\begin{enumerate}
    \item At the top of \texttt{shadeWaterSurface()} (right after \texttt{screenUV} is computed,
    \texttt{water\_surface.glsl:601--604}), read \texttt{float solidD = textureLod(solidSceneDepthTex, screenUV,
    0.0).r;} and \texttt{discard;} when \texttt{solidD < gl\_FragCoord.z - kDepthEps} with a small positive epsilon
    (the same convention the composite uses, \texttt{postprocess.frag:148}, tightened so coplanar shorelines are not
    rejected).
    \item Guard the fetch so it cannot fire when the solid depth is not valid for this pass: add a UBO/pass flag (or
    check \texttt{ubo.passParams} / the existing capture-mode conventions) so the 360-capture and any future
    water-in-main variant skip the test (there the hardware depth test already does it).
    \item Keep the composite test (it also covers vegetation and remains correct if the shader-level test is disabled).
    \item Fix the comment at \texttt{WaterRenderer.cpp:1250} to state that occlusion is resolved by the shader
    early-out plus the composite test, not by the forward pass.
    \item Measure with a camera placed behind a ridge overlooking a lake: compare GPU timestamps 14/15 with and without
    the change.
\end{enumerate}

\textbf{Files to modify:}
\begin{itemize}
    \item \texttt{shaders/includes/water\_surface.glsl} (early-out at \texttt{:601--604})
    \item \texttt{vulkan/renderer/WaterRenderer.cpp} (comment at \texttt{:1250}; optional pass flag)
    \item \texttt{vulkan/ubo/WaterRenderUBO.hpp} (only if a new flag channel is needed)
\end{itemize}

\textbf{Validation:} RenderDoc shows the discard in the water pipeline and a large drop in water-shader invocations for
occluded views; screenshots with the camera above open water are pixel-identical (no shoreline is rejected --- verify at
grazing angles and at the waterline); no validation errors from the depth fetch (the image is already in
\texttt{SHADER\_READ\_ONLY}); GPU-assisted validation clean.
\end{tcolorbox}

% =====================================================================
\section{High-Severity Issues}
% =====================================================================

\subsection{H6: The Water Vertex Stage Is 8.5k Instructions per Vertex and Runs Twice per Frame}

\textbf{Location:} \texttt{shaders/main.vert:134--136} (\texttt{waterDisplaceWaterVertex} in the WATER\_NO\_TESS path;
module \texttt{main\_water\_no\_tess.vert.spv} = 8,464 instructions, 513 PCG sites),
\texttt{shaders/includes/water\_tese.glsl:113--152} (the shared core), \texttt{:345} (TES call site; module
\texttt{main\_water.tese.spv} = 8,604 instructions); \texttt{vulkan/renderer/WaterBackFaceRenderer.cpp:243--254}
(back-face pass re-runs the same vertex path with front-face culling).

\begin{itemize}
    \item With tessellation off --- the default --- the water vertex shader still calls
    \texttt{waterDisplaceWaterVertex}, which evaluates the full wave field: 4 FBM sites $\times$ 4 octaves =
    \textbf{16 four-dimensional Perlin evaluations $\approx$ 2,048 PCG hashes per vertex}. The module is 8,464
    instructions against \textbf{137} for the solid vertex shader.
    \item It runs \textbf{twice} per frame for every water vertex: once in the geometry pass and once in the back-face
    depth pass, which draws the same mesh with \texttt{VK\_CULL\_MODE\_FRONT\_BIT}.
    \item Most of that work cannot be represented by the geometry anyway. The base mesh is a chunk-resolution isosurface
    whose vertices are metres apart, while the finest octave of the default spectrum has a 6.25\,cm wavelength: the top
    octaves are pure aliasing on geometry. Band-limiting the per-vertex evaluation is both a correctness improvement
    (no aliased silhouette) and a $\sim$2$\times$ saving.
    \item The back-face pass needs a \emph{depth proxy}, not a displaced surface: it exists only so the fragment stage
    can measure a water column. Displacing its vertices with the wave field is work whose only effect is to add noise to
    that measurement.
    \item The vertex stage additionally performs 2--6 depth-texture fetches per vertex (base depth, shore-gradient pair,
    plus the displaced-UV refinement at \texttt{main.vert:110--133}), and the whole thing is gated only by
    \texttt{wp.waveToggles.x > 0.5}, never by distance.
\end{itemize}

\textbf{Technique:} band-limit the per-vertex field (drop octaves finer than the local vertex spacing), gate the
displacement by distance, and give the back-face pass an undisplaced vertex variant.

\begin{tcolorbox}[promptbox={Prompt: Band-Limit the Vertex Wave Evaluation and Stop Displacing Back Faces}]
\textbf{Goal:} The water vertex stage costs a few octaves instead of 16 Perlin4D per vertex, and the back-face pass pays
no wave evaluation at all.

\textbf{Steps:}
\begin{enumerate}
    \item Add an explicit octave budget to \texttt{waterDisplaceWaterVertex}
    (\texttt{water\_tese.glsl:113}) and clamp \texttt{octavesN} inside \texttt{waterWaveField} with it. Derive it from
    the local vertex spacing: pass the patch edge length (TES) or use \texttt{fwidth}-free approximations such as the
    camera-distance band already used for tessellation (VS), and drop every octave whose wavelength is below $\sim$2
    vertex spacings.
    \item Gate the displacement by distance: beyond the distance where the displacement subtends less than a pixel
    (or beyond the existing 200--1000\,m detail band, consistent with C4), skip the field call entirely and emit the
    base vertex. Keep \texttt{fragWaterDepth}/\texttt{fragShoreDir} --- the depth measurement still matters for the
    region tint --- but skip the field.
    \item Add an undisplaced back-face vertex variant: a \texttt{WATER\_BACKFACE} define (or a dedicated
    \texttt{water\_backface} vertex module) that emits the base position and forwards
    \texttt{fragWaterDepth}/\texttt{fragShoreDir} without calling \texttt{waterDisplaceWaterVertex}. Wire it into
    \texttt{WaterBackFaceRenderer}'s pipeline family selection. Document that the back-face depth is a volume proxy and
    that the wave displacement perturbs it by less than the 5\,cm validity threshold
    (\texttt{water\_surface.glsl:638}).
    \item Share the per-vertex result with the fragment stage where possible (M11) so the two evaluations stop
    duplicating each other.
    \item Re-measure: expect \texttt{main\_water\_no\_tess.vert} / \texttt{main\_water.tese} to drop from $\sim$8.5k
    instructions to a fraction of that, and the PCG-site count from 513 to $\sim$128--256.
\end{enumerate}

\textbf{Files to modify:}
\begin{itemize}
    \item \texttt{shaders/main.vert} (no-tess path), \texttt{shaders/includes/water\_tese.glsl} (shared core, budget)
    \item \texttt{vulkan/renderer/WaterBackFaceRenderer.cpp/hpp} (undisplaced variant)
    \item \texttt{Makefile} (new vertex variant rule)
\end{itemize}

\textbf{Validation:} Wireframe and shoreline screenshots are unchanged (the silhouette must not visibly change where the
displacement is retained); screenshots of distant water show no popping; the back-face depth debug view
(\texttt{DEBUG\_MODE\_THICKNESS}) is within tolerance; validation clean (the new variant's vertex/fragment interface must
match the back-face fragment stage), GPU-assisted validation clean.
\end{tcolorbox}

\subsection{H7: The TCS Evaluates the Full Wave Field Three Times per Patch to Nudge Tessellation by 30\%}

\textbf{Location:} \texttt{shaders/includes/water\_tesc.glsl:84--106}; module \texttt{main\_water.tesc.spv} = 7,682
instructions, 513 PCG sites; parameter \texttt{utils/WaterParams.hpp:91} (\texttt{tessNoiseInfluence = 0.3}).

\begin{itemize}
    \item For each of the three patch edges the TCS evaluates \texttt{waterWaveDisplacement} at the edge midpoint
    (\texttt{water\_tesc.glsl:98--100}) --- a full field evaluation (4 FBM sites $\times$ 4 octaves = 16 Perlin4D) ---
    purely to multiply the distance-derived tessellation level by
    \texttt{1.0 + noiseInf * (noiseVal - 0.5)}, i.e. by at most $\pm30\%$ at the default.
    \item That is \textbf{48 four-dimensional Perlin evaluations per input triangle}, on top of the TES evaluation for
    every generated vertex, for a quantity whose only job is to bias triangle density by a few percent.
    \item The noise must stay deterministic per shared edge (so adjacent patches agree and no cracks appear), which
    rules out a screen-space or per-patch shortcut --- but not a cheaper field: the tessellation level does not need the
    chop, the calm mask, the second train, or four octaves.
    \item This cost only exists when \texttt{settings.tessellationEnabled} is true (default false), but it is the
    dominant per-patch cost of the maximum-quality path.
\end{itemize}

\textbf{Technique:} add a dedicated single-train, one-or-two-octave tessellation probe that keeps the edge-midpoint
determinism, or reuse the TES-driven density and drop the noise term when \texttt{tessNoiseInfluence} is small.

\begin{tcolorbox}[promptbox={Prompt: Replace the TCS Wave Probe with a Single-Train, Low-Octave Estimate}]
\textbf{Goal:} The noise-adaptive tessellation modulation costs $\leq$4 four-dimensional Perlin evaluations per patch
instead of 48.

\textbf{Steps:}
\begin{enumerate}
    \item Add \texttt{float waterWaveTessProbe(vec3 pos, float time, WaterParamsGPU wp)} to
    \texttt{water\_noise.glsl}: one \texttt{waterRidgedFbmGrad} on the \emph{primary} train only, with
    \texttt{octaves = min(2, octavesN)}, no chop, no mask, no foam. Keep the same \texttt{shoreDir} fallback and the same
    deep-water (\texttt{depth = -1}) convention so the probe tracks the real surface.
    \item Use it at \texttt{water\_tesc.glsl:98--100}. Keep evaluating at the edge midpoint so the crack-free property
    (identical value for both patches sharing an edge) is preserved exactly; add a comment restating why.
    \item Skip the probe entirely when \texttt{noiseInf <= 0.0} or \texttt{wp.waveToggles.x < 0.5} (already partly
    handled at \texttt{:93}) --- and consider folding \texttt{tessNoiseInfluence} into the distance term when it is
    below a threshold, documenting the visual trade-off.
    \item Re-measure: expect \texttt{main\_water.tesc} PCG sites to drop from 513 to $\sim$64--128.
\end{enumerate}

\textbf{Files to modify:}
\begin{itemize}
    \item \texttt{shaders/includes/water\_noise.glsl} (new probe), \texttt{shaders/includes/water\_tesc.glsl}
    (call site \texttt{:98--100})
\end{itemize}

\textbf{Validation:} Tessellation debug heat view shows the same density distribution within a small tolerance; no cracks
appear at patch borders (fly along a shoreline at maximum quality); PCG-site count drops; validation clean.
\end{tcolorbox}

\subsection{H8: The Water Color Target Is RGBA32F at Full Resolution}

\textbf{Location:} \texttt{vulkan/renderer/WaterRenderer.cpp:340} (allocation, \texttt{VK\_FORMAT\_R32G32B32A32\_SFLOAT}),
\texttt{:772, 861--869} (pipeline attachment formats), \texttt{:1171, 1308, 1358} (transitions),
\texttt{shaders/postprocess.frag:98} (composite fetch).

\begin{itemize}
    \item The water color target is 128 bits per pixel at full resolution: 16 bytes written per water pixel, 16 bytes
    read per screen pixel in the composite, plus a full-res clear and two layout transitions per frame. Nothing in the
    shader needs float32: the color is a lit, tone-mapped-range value with a handful of \texttt{pow}-shaped highlights
    and an alpha in $[0,1]$.
    \item RGBA16F halves both the write and the read bandwidth and the clear cost, at the same time reducing the
    pressure that makes the (already huge) water fragment shader memory-bound at 4K.
    \item The body/column attachments are already 16F/16F and were made conditional by report~19's H4; the color target
    is the last full-precision target in the water path.
\end{itemize}

\textbf{Technique:} switch the water color target and its pipeline attachment format to
\texttt{VK\_FORMAT\_R16G16B16A16\_SFLOAT}; keep the dummy water color view format in sync.

\begin{tcolorbox}[promptbox={Prompt: Move the Water Color Target to RGBA16F}]
\textbf{Goal:} Half the water-target bandwidth with no visible change.

\textbf{Steps:}
\begin{enumerate}
    \item Replace \texttt{VK\_FORMAT\_R32G32B32A32\_SFLOAT} with \texttt{VK\_FORMAT\_R16G16B16A16\_SFLOAT} for the water
    color image (\texttt{WaterRenderer.cpp:340}), its transitions (\texttt{:348, 487, 527}), the pipeline rendering
    formats (\texttt{:861}) and the no-body variant (\texttt{:892}). Confirm the dummy water color image
    (\texttt{:283--290}, already RGBA8) stays consistent with whatever the composite samples.
    \item Update \texttt{renderMainTargets}' format handling if the water-in-main path shares this format constant.
    \item Check the water-in-main path: there the attachment is the \emph{solid} color target, so no change is needed
    there --- but confirm no code assumes the offscreen water format equals the solid one.
    \item Verify no debug view needs the extra range (the displacement/thickness debug views write normalised values;
    the compose view writes clamped fractions).
\end{enumerate}

\textbf{Files to modify:}
\begin{itemize}
    \item \texttt{vulkan/renderer/WaterRenderer.cpp/hpp} (allocation, transitions, pipeline formats, dummy)
\end{itemize}

\textbf{Validation:} Screenshots are within 8-bit tolerance at both default and maximum quality (compare with
\texttt{bin/screenshot.png} before/after); specular highlights and caustics do not clip; VRAM budget drops by
$3\times$ (width $\times$ height $\times$ 8 bytes); validation clean (format/layout pairs must match everywhere).
\end{tcolorbox}

\subsection{H9: A Separate Full-Resolution Water Pass Survives Even Though the Water-in-Main Path Exists}

\textbf{Location:} \texttt{MyApp.cpp:2583--2616} (\texttt{settings.waterInMainPass} selects
\texttt{renderMainTargets} vs \texttt{renderPass}), \texttt{utils/Settings.hpp:133} (default \texttt{false}),
\texttt{widgets/SettingsWidget.cpp:142} (the toggle), \texttt{vulkan/renderer/WaterRenderer.cpp:997--1120}
(the offscreen pass), \texttt{shaders/postprocess.frag:98, 146--151} (composite water fetches).

\begin{itemize}
    \item In the default path water is rendered into its own full-resolution color $+$ depth pair, transitioned twice
    per frame, cleared every frame, and then fetched twice more by the composite (water color at \texttt{:98}, water
    geometry depth at \texttt{:147}) for every screen pixel.
    \item The Phase-1 water-in-main path (\texttt{renderMainTargets}, the alpha-blended water pipeline) already exists
    and is user-selectable. Drawing water into the main targets removes the offscreen pair, its transitions and clears,
    and --- decisively --- gives the water draw \textbf{hardware early-Z against the solid depth buffer}, which is the
    structural fix for C5 rather than a shader-level discard.
    \item It is not free: the in-trace screen lookups then bind the previous frame's solid/vegetation views (1-frame
    latency, as documented at \texttt{MyApp.cpp:2588--2592}), and the composite's separate-water branch is skipped.
    \item Independently, a \textbf{water render scale} (e.g.\ 0.5) is a standard technique for translucent layers: the
    water surface is smooth and low-frequency compared with the scene, and a bilinear upsample of color plus a
    closest-depth upsample of the water depth is visually acceptable for the composite's occlusion test.
\end{itemize}

\textbf{Technique:} finish Phase~1b (make water-in-main the default, with the 1-frame latency documented and A/B'd) and,
for the offscreen path, add a resolution scale.

\begin{tcolorbox}[promptbox={Prompt: Make Water-in-Main the Default and Add a Water Resolution Scale}]
\textbf{Goal:} Either water shares the main targets (no offscreen pair, hardware early-Z), or it is shaded at a
fraction of the resolution --- in both cases the per-frame water cost drops substantially.

\textbf{Steps:}
\begin{enumerate}
    \item \textbf{Phase 1b:} flip \texttt{Settings::waterInMainPass} to \texttt{true}, run the full smoke test, and
    quantify the 1-frame latency on the in-trace lookups (camera motion near the waterline and under fast rotation).
    Fix any ordering issue rather than reverting the default: the water task already waits \texttt{tlSolid}.
    \item Verify the composite's dummy-water path
    (\texttt{MyApp.cpp:3083--3092}) leaves a fully transparent water view so no ghosting appears, and that the brush
    overlay and back-face depth paths still work on the in-main route.
    \item \textbf{Resolution scale:} add \texttt{Settings::waterRenderScale} (1.0 / 0.5), pass
    \texttt{width * scale} to \texttt{WaterRenderer::createRenderTargets} and
    \texttt{WaterBackFaceRenderer::createRenderTargets}, and keep the viewport/scissor in device pixels (they already
    derive from \texttt{renderWidth/renderHeight}).
    \item In the composite, sample the water color with a bilinear fetch (it already is one) and sample the water
    geometry depth with an explicit \texttt{textureLod(..., 0.0)}; for the occlusion test take the \emph{minimum} of the
    2$\times$2 depth taps (closest-surface upsample) so terrain edges in front of water are never punched through.
    Document the half-pixel waterline tolerance this introduces.
    \item Keep the full-resolution path as the reference for A/B and for screenshots.
\end{enumerate}

\textbf{Files to modify:}
\begin{itemize}
    \item \texttt{utils/Settings.hpp}, \texttt{MyApp.cpp} (target sizes, in-main default), \texttt{shaders/postprocess.frag}
    \item \texttt{vulkan/renderer/WaterRenderer.cpp/hpp}, \texttt{vulkan/renderer/WaterBackFaceRenderer.cpp/hpp}
    \item \texttt{vulkan/renderer/SceneRenderer.cpp} (target creation at \texttt{:181--186, 487--489})
\end{itemize}

\textbf{Validation:} Water-in-main: screenshots match the offscreen path within tolerance on static views; no water
ghosting; brush liquid still composites. Half-res: the waterline shows no terrain bleeding through (closest-depth
upsample), and reflection/specular detail loss is documented and accepted; VRAM drops by $\sim$75\% of the water
targets; validation clean including the new image sizes and layout transitions.
\end{tcolorbox}

% =====================================================================
\section{Medium-Severity Issues}
% =====================================================================

\subsection{M10: Specular and Glitter Run 10 Perlin Evaluations per Pixel With No Lobe Early-Out}

\textbf{Location:} \texttt{shaders/includes/water\_surface.glsl:1174--1190}; defaults
\texttt{utils/WaterParams.hpp:57,59} (\texttt{specularIntensity = 2.0}, \texttt{glitterIntensity = 1.5}).

\begin{itemize}
    \item The block evaluates three FBM chains: the specular perturbation (4 octaves), the glitter noise
    (\texttt{noiseOctaves - 2} = 2) and the glitter threshold (4) --- \textbf{10 four-dimensional Perlin evaluations},
    measured at 384 of the module's 3,073 PCG sites.
    \item Both lobes are multiplied by powers of \texttt{specAngle}: \texttt{pow(specAngle, specularPowerParam)} with
    \texttt{specularPower = 128}, and \texttt{pow(specAngle, 32.0)} for the glitter. Outside the sun lobe those powers
    are numerically zero, so the noise is computed to be multiplied by zero for the overwhelming majority of pixels.
    \item Report~19 gated this block on \texttt{wp.waveToggles.x > 0.5}, which helps calm layers, but the lobe gate is
    still missing: with waves on, every water pixel pays it.
\end{itemize}

\textbf{Technique:} hoist the cheap lobe term and early-out on it; scale the glitter noise cost by the lobe too.

\begin{tcolorbox}[promptbox={Prompt: Gate the Specular/Glitter Noise on the Highlight Lobe}]
\textbf{Goal:} Zero noise cost for the $\sim$99\% of water pixels outside the sun highlight.

\textbf{Steps:}
\begin{enumerate}
    \item Compute \texttt{specAngle} (already available at \texttt{water\_surface.glsl:1172}) and the lobe powers
    \emph{before} the FBMs. Skip the specular FBM when \texttt{pow(specAngle, specularPowerParam) * specularIntensity}
    is below a visible threshold (e.g.\ $10^{-4}$), and skip both glitter FBMs when \texttt{pow(specAngle, 32.0)} is
    below theirs.
    \item Merge the two glitter chains: the threshold FBM (4 octaves) exists only to jitter a cutoff. Reuse the glitter
    noise value itself with an offset/threshold remap, halving the glitter cost.
    \item Apply the distance budget from C4 to this block as well (glitter at 1\,km is invisible).
    \item Keep the analytic highlight for the waves-off branch (\texttt{:1191--1194}) unchanged --- it is already free.
\end{enumerate}

\textbf{Files to modify:}
\begin{itemize}
    \item \texttt{shaders/includes/water\_surface.glsl} (\texttt{:1174--1190})
\end{itemize}

\textbf{Validation:} Screenshots with the sun in view are pixel-identical (the lobe is unchanged where it is visible);
screenshots with the sun behind the camera show no change because the lobe was already $\sim$0; a profiling counter or
debug view confirms the blocks are skipped outside the lobe; validation clean.
\end{tcolorbox}

\subsection{M11: The Vertex Stage and the Fragment Stage Evaluate the Same Field at the Same Position}

\textbf{Location:} \texttt{shaders/includes/water\_tese.glsl:113--152} / \texttt{shaders/main.vert:134--136}
(\texttt{waterDisplaceWaterVertex} $\rightarrow$ \texttt{waterWaveSample}), and
\texttt{shaders/includes/water\_surface.glsl:665--667} (the fragment-stage \texttt{waterWaveField} at
\texttt{fragBasePos.xyz} with the same \texttt{animTime}, \texttt{fragWaterDepth} and \texttt{fragBasePos.w}).

\begin{itemize}
    \item Both stages evaluate the identical field at the identical base position with identical parameters. The vertex
    stage uses the result to displace geometry and to derive an interpolated normal; the fragment stage re-evaluates it
    per pixel for the analytic normal, because interpolation cannot carry the fine octaves.
    \item The duplication is not pure waste --- the per-pixel evaluation is what gives the surface its detail --- but the
    \emph{coarse} part of the spectrum is computed twice at two different rates for the same answer.
    \item With the octave budget from C4/H6 in place, the natural split becomes available: the vertex stage evaluates the
    low octaves (which the geometry can represent) and forwards them; the fragment stage evaluates only the octaves above
    that cutoff.
\end{itemize}

\textbf{Technique:} split the spectrum by evaluation rate once the octave budget exists; until then, at minimum, forward
the vertex-stage height/gradient so the fragment stage can skip the octaves the vertex stage already resolved.

\begin{tcolorbox}[promptbox={Prompt: Split the Wave Spectrum Between the Vertex and Fragment Rates}]
\textbf{Goal:} Each octave of the wave field is evaluated exactly once per frame at the rate that can represent it.

\textbf{Steps:}
\begin{enumerate}
    \item Give \texttt{waterWaveField} an explicit \texttt{octLo}/\texttt{octHi} range (in addition to the C4 octave
    budget) so a caller can request ``octaves 0--1 only'' or ``octaves 2--3 only''.
    \item In the vertex stage, request the low band and forward its height and analytic gradient through two extra
    varyings (or pack the gradient into the existing \texttt{fragDebug}/\texttt{fragBasePos} channels if locations are
    scarce --- prefer new locations and update \texttt{vulkan/includes/locations.hpp}).
    \item In the fragment stage, request only the high band and add the forwarded low-band gradient to it before
    building the normal. The result must be identical to evaluating the full spectrum per pixel.
    \item Verify the shore-zone envelope (\texttt{env}, \texttt{heightTaper}, \texttt{mask}) is applied consistently:
    those are per-point scalars; compute them once (cheaply) in both stages or forward them too.
    \item If the added varyings cost more than the saved octaves on your target (measure!), keep the current duplication
    and instead rely on H6's band-limiting alone.
\end{enumerate}

\textbf{Files to modify:}
\begin{itemize}
    \item \texttt{shaders/includes/water\_noise.glsl} (octave range), \texttt{shaders/includes/water\_tese.glsl},
    \texttt{shaders/main.vert}, \texttt{shaders/includes/water\_surface.glsl}
    \item \texttt{vulkan/includes/locations.hpp} (new varyings, if taken)
\end{itemize}

\textbf{Validation:} \texttt{DEBUG\_MODE\_SHADING\_NORMAL} and \texttt{DEBUG\_MODE\_DISPLACEMENT} are identical to the
baseline at near and far range; no discontinuity appears across chunk borders where the vertex density changes (this is
the main risk --- test at a LoD transition); validation clean (varying interface must match in every water pipeline
including the back-face and wireframe ones).
\end{tcolorbox}

\subsection{M12: The Composite Fetches the Water Geometry Depth for Every Screen Pixel}

\textbf{Location:} \texttt{shaders/postprocess.frag:98} (water color fetch), \texttt{:146--151} (the occlusion block).

\begin{itemize}
    \item The occlusion block fetches \texttt{waterGeomDepthTex} unconditionally for every pixel of the frame --- sky,
    terrain, vegetation, brush --- even though \texttt{waterAlpha == 0} makes the test a no-op there. Two full-res
    fetches (color $+$ depth) are therefore paid on every non-water pixel.
    \item The composite already runs at full resolution over the whole screen, so this is a bandwidth and latency cost
    proportional to the display resolution.
\end{itemize}

\textbf{Technique:} guard both fetches and the occlusion test on \texttt{waterAlpha > 0.0}.

\begin{tcolorbox}[promptbox={Prompt: Skip the Water Composite Work Where There Is No Water}]
\textbf{Goal:} Zero water fetches for pixels the water pass did not cover.

\textbf{Steps:}
\begin{enumerate}
    \item Move the \texttt{waterGeomDepthTex} fetch and the occlusion test inside \texttt{if (waterAlpha > 0.0)}
    (\texttt{postprocess.frag:146--151}). Note that \texttt{waterAlpha} comes from the color fetch, which must stay
    unconditional (it is the coverage signal) --- but the depth fetch, the \texttt{waterGeomDepth} comparison and the
    alpha zeroing do not.
    \item Fold the blur block (\texttt{:112--140}) and the occlusion block into one guarded region, preserving the
    current order (blur first, then occlusion).
    \item Keep the \texttt{waterBlurEnabled} gate from report~19's H4 intact.
\end{enumerate}

\textbf{Files to modify:}
\begin{itemize}
    \item \texttt{shaders/postprocess.frag}
\end{itemize}

\textbf{Validation:} Screenshots pixel-identical on water and non-water regions; a camera with no water on screen shows
zero water-depth fetches in a RenderDoc capture; validation clean.
\end{tcolorbox}

\subsection{M13: Water That Is Resident but Entirely Frustum-Culled Still Pays the Full Pass}

\textbf{Location:} \texttt{MyApp.cpp:2214--2245} (the empty-water fast path), \texttt{:2224--2228}
(\texttt{waterMeshCount} is the \emph{resident} mesh count).

\begin{itemize}
    \item Report~19's M8 fast path fires only when there are \emph{zero resident} water meshes. As soon as one water
    chunk is loaded anywhere in the scene, every frame performs the full sequence even when the camera looks away from
    all water: descriptor updates, the back-face path, three clears $+$ a depth clear, an indirect draw with count 0, and
    the layout transitions.
    \item The M8 comment (\texttt{MyApp.cpp:2216--2223}) explains why the GPU visible-count readback was rejected: it is
    copied asynchronously and can lag several frames, so a stale zero cleared the targets while the current cull still
    had visible chunks, producing flicker while chunks streamed in.
    \item The fix is not to trust a single stale read but to require \emph{persistence}: a value that has been zero for
    $k$ consecutive frames is safe, because the fast path only ever clears the targets to alpha 0 (so a wrong guess costs
    one frame of missing water, and a persistent zero means the camera has been looking away for $k$ frames).
\end{itemize}

\textbf{Technique:} extend the fast path with an N-frame hysteresis on the GPU visible count, keeping the clear-to-alpha-0
safety property.

\begin{tcolorbox}[promptbox={Prompt: Add Hysteresis to the Empty-Water Fast Path}]
\textbf{Goal:} The clear-only path also fires when the water is resident but has been invisible for several frames.

\textbf{Steps:}
\begin{enumerate}
    \item Read the water indirect renderer's visible count (the same source
    \texttt{SceneRenderer::readVisibleCount} exposes) and keep a small per-frame history or a counter
    \texttt{waterInvisibleFrames}.
    \item Take the fast path when \texttt{waterInvisibleFrames >= k} (start with $k = 3$) \emph{and} no brush liquid is
    pending. Reset the counter to 0 on any frame with a nonzero count. Because the fast path clears the water color to
    $(0,0,0,0)$, a mistaken skip shows one frame of missing water at worst --- and with $k \ge 2$ a streamed-in chunk
    cannot be missed, since the cull runs before the draw in the same frame.
    \item Keep the existing zero-resident-meshes path unconditional (it needs no hysteresis).
    \item Do not apply this to the water-in-main variant (the composite reads the main targets there).
\end{enumerate}

\textbf{Files to modify:}
\begin{itemize}
    \item \texttt{MyApp.cpp} (water task, \texttt{:2214--2245}), \texttt{vulkan/renderer/SceneRenderer.hpp/cpp}
    (expose the visible count if it is not already reachable)
\end{itemize}

\textbf{Validation:} Camera-away captures show no water raster passes after $k$ frames; spinning the camera back toward
water shows no flicker or one-frame flash; chunk streaming does not produce flicker (reproduce the case the M8 comment
describes); validation clean.
\end{tcolorbox}

% =====================================================================
\section{Low-Severity / Hygiene Issues}
% =====================================================================

\subsection{L14: One Noise Spectrum Serves Five Consumers With Wildly Different Band Requirements}

\textbf{Location:} \texttt{utils/WaterParams.hpp:37--44} (\texttt{noiseOctaves = 4},
\texttt{noisePersistence = 0.5}, \texttt{noiseLacunarity = 4.0}, \texttt{noisePeriod = 4.0}); consumers
\texttt{shaders/includes/water\_noise.glsl:225} (chop), \texttt{:239} (calm mask), \texttt{:271,273} (ridged trains),
\texttt{:361} (foam texture), \texttt{:400--410} (refraction offset); packing
\texttt{vulkan/ubo/WaterParamsGPU.hpp:7}.

\begin{itemize}
    \item All five consumers read the same \texttt{params2.z} (octaves), \texttt{params2.w} (persistence) and
    \texttt{params3.y} (lacunarity). A user who wants fine foam detail therefore buys four octaves of ridged swell they
    cannot see, and a user who wants broad swell buys four octaves of foam texture.
    \item \texttt{noiseLacunarity = 4.0} is unusual (2.0--2.5 is the conventional range). With four octaves it spans a
    $64\times$ frequency range, so the top octaves are sub-pixel at almost any distance: paid in full, visible not at
    all, and a source of the shimmer C4's footprint band-limiting would remove.
    \item Splitting the spectrum is a small packing change with a large multiplier effect once C1--C4 land, because each
    consumer can then be band-limited independently.
\end{itemize}

\textbf{Technique:} give the refraction offset, the chop/mask and the foam their own octave/lacunarity fields (or derive
them from the existing ones with documented per-consumer clamps), and lower the lacunarity default.

\begin{tcolorbox}[promptbox={Prompt: Give Each Noise Consumer Its Own Band}]
\textbf{Goal:} No consumer pays for octaves it cannot show; the default lacunarity stops producing sub-pixel octaves.

\textbf{Steps:}
\begin{enumerate}
    \item Lower \texttt{noiseLacunarity} to $\sim$2.2 and re-tune \texttt{noisePeriod} so the \emph{base} feature size is
    unchanged; A/B the look and document the new defaults.
    \item Add per-purpose octave counts: reuse the free channels of \texttt{WaterParamsGPU} (or extend the struct) with
    \texttt{refractionOctaves}, \texttt{foamOctaves} and \texttt{detailOctaves}, defaulting each to a value derived from
    the current \texttt{noiseOctaves} so the shipped look is unchanged.
    \item Thread them through \texttt{WaterRenderer}'s packing into \texttt{water\_noise.glsl} and surface them in
    \texttt{widgets/WaterWidget.cpp} under the existing noise section (optional but cheap).
    \item Clamp every consumer's octave count against the footprint-based budget from C4 so a user cannot configure
    their way back to sub-pixel octaves.
\end{enumerate}

\textbf{Files to modify:}
\begin{itemize}
    \item \texttt{utils/WaterParams.hpp}, \texttt{vulkan/ubo/WaterParamsGPU.hpp},
    \texttt{vulkan/renderer/WaterRenderer.cpp} (packing)
    \item \texttt{shaders/includes/water\_noise.glsl}, \texttt{shaders/includes/water\_surface.glsl}
    \item \texttt{widgets/WaterWidget.cpp} (optional UI)
\end{itemize}

\textbf{Validation:} Default screenshots unchanged after the lacunarity re-tune (A/B required); the new octave fields
default to the current look; no validation errors from the UBO layout change (check the SSBO stride and the
\texttt{waterParams.length()} indexing used at \texttt{water\_surface.glsl:559}).
\end{tcolorbox}

% =====================================================================
\section{Implementation Roadmap and Expected Deltas}
% =====================================================================

\begin{table}[h]
\centering
\small
\begin{tabular}{@{}llp{6.6cm}@{}}
\toprule
\textbf{Order} & \textbf{Issue} & \textbf{Expected delta} \\
\midrule
1 & C1 & $-62.5\%$ hash ALU everywhere; 46.3k $\rightarrow$ 20.5k instructions in the water fragment module. \\
2 & C2 & Removes 13 Perlin4D/pixel --- i.e. essentially all of the default preset's procedural water cost --- plus a distance fade. \\
3 & C4 & Removes the field/spec/glitter/caustic cost on all water beyond the detail band (usually the majority of water pixels). \\
4 & C5 & Removes the full shader for every terrain-occluded water pixel. \\
5 & C3 & 32 $\rightarrow$ 4--8 Perlin4D/pixel for caustics; the largest single term in a wavy frame. \\
6 & H6 & 16 $\rightarrow$ few Perlin4D per water vertex, and none at all in the back-face pass. \\
7 & H7 & 48 $\rightarrow$ $\le$4 Perlin4D per tessellated patch. \\
8 & H8 & Halves the water-target write, read and clear bandwidth. \\
9 & H9 & Removes the offscreen pair entirely (water-in-main) or quarters the shaded pixels (half-res). \\
10 & M10 & Removes 10 Perlin4D/pixel outside the sun lobe. \\
11 & M11 & Removes the duplicated coarse spectrum between the two stages. \\
12 & M12 & Removes two full-res fetches per non-water pixel. \\
13 & M13 & Removes the pass, clears and transitions when water is resident but invisible. \\
14 & L14 & Stops every consumer paying for octaves it cannot resolve. \\
\bottomrule
\end{tabular}
\end{table}

\textbf{Measurement protocol.} Keep the \S\ref{sec:method} ablation harness in the tree (a scratch copy of
\texttt{shaders/} plus the \texttt{glslc} invocation) so each landing can be re-measured: PCG-site count, SPIR-V
instruction count, and the op-code mix. Cross-check with the existing GPU timestamps 14/15 (back-face $+$ water
geometry, \texttt{MyApp.cpp:2582, 2617}) and with RenderDoc (bound pipeline stages, texture-fetch counters,
\texttt{indirect.comp} dispatches). Keep the RT-enabled presets as the regression baseline: no issue above is allowed to
change their output. \texttt{AGENTS.md} rules apply to every prompt: runtime feature detection, Synchronization2,
reverse-order destruction, no legacy render passes, and validation-clean runs
(\texttt{VULKAN\_GPU\_ASSISTED=1 make run-debug}).

\begin{table}[h]
\centering
\small
\begin{tabular}{@{}lrrr@{}}
\toprule
\textbf{Water fragment module / frame} & \textbf{Today} & \textbf{After C1--C4} & \textbf{After all} \\
\midrule
SPIR-V instructions (static) & 46,320 & $\sim$20.5k & $\sim$12k \\
Inlined PCG hash sites & 3,073 & $\sim$1,153 & $\sim$600 \\
Perlin4D per pixel, default preset & 13 & 0--2 & 0--2 \\
Perlin4D per pixel, waves $+$ caustics & 75 & $\sim$10 near, 0 far & $\sim$6 near, 0 far \\
Perlin4D per water vertex ($\times 2$ passes) & 16 & 4 & 0--2 (front only) \\
Water target bytes per pixel (write $+$ read) & 32 & 32 & 0--16 \\
\bottomrule
\end{tabular}
\caption{Qualitative budget; exact counts depend on the final variant choices and on the ABI of L14's spectrum split.}
\end{table}

\end{document}
